\id{МРНТИ 68.75.13}{}

\begin{header}
\swa{А.М.~Шильманова, Г.А.~Ибраева, А.С.~Шайнуров, Г.К.~Иляшова, С.А.~Юсупова, С.Ғ.~Тұңғышбай}{СИСТЕМНЫЕ ФАКТОРЫ ФОРМИРОВАНИЯ ПРОДОВОЛЬСТВЕННОЙ БЕЗОПАСНОСТИ АГРОПРОДОВОЛЬСТВЕННОЙ СИСТЕМЫ}

\tsp{1}А.М.~Шильманова\envelope,
\tsp{1}Г.А.~Ибраева,
\tsp{1}А.С.~Шайнуров,
\tsp{2}Г.К.~Иляшова,
\tsp{3}С.А.~Юсупова,
\tsp{1}С.Ғ.~Тұңғышбай
\end{header}

\begin{affil}
\tsp{1}Кызылординский университет имени Коркыт Ата, Кызылорда, Казахстан,

\tsp{2}Казахский национальный университет им. аль-Фараби, Алматы, Казахстан,

\tsp{3}Казахская национальная академия хореографии, Астана, Казахстан

\corrauthor{Корреспондент-автор: altynajsilmanova1@gmail.com}
\end{affil}

В условиях нарастающей глобальной нестабильности и усиления внешних и
внутренних рисков продовольственная безопасность все в большей степени
приобретает системный характер и требует переосмысления традиционных
подходов к ее анализу и обеспечению. Современные агропродовольственные
системы функционируют в условиях высокой неопределенности, что
актуализирует необходимость рассмотрения продовольственной безопасности
не только через показатели производства и доступности продовольствия, но
и через характеристики устойчивости, согласованности и воспроизводимости
взаимодействий между ключевыми звеньями системы.

Целью исследования является выявление и анализ системных факторов
формирования продовольственной безопасности агропродовольственной
системы, а также определение их роли в обеспечении устойчивости
продовольственного обеспечения в условиях внешних и внутренних шоков.
Методологическую основу исследования составляют системный,
институциональный и сравнительный подходы, позволяющие рассматривать
продовольственную безопасность как результат функционирования
агропродовольственной системы в целом. Информационной базой послужили
материалы международных организаций, научные публикации и официальные
статистические данные, включая источники ФАО, ОЭСР и национальные
аналитические материалы.

В результате исследования показано, что в Республике Казахстан модель
обеспечения продовольственной безопасности в значительной степени
опирается на финансово-инструментальные меры государственной поддержки,
ориентированные на отдельных хозяйствующих субъектов. При этом механизмы
координации, кооперации и устойчивой интеграции участников
агропродовольственных цепочек развиты в ограниченной степени и не
выполняют системообразующей функции. Сделан вывод о том, что
устойчивость продовольственного обеспечения в долгосрочной перспективе
определяется не столько масштабами бюджетной поддержки, сколько
качеством институциональной среды и воспроизводимостью взаимодействий
между звеньями агропродовольственной системы.

{\bfseries Ключевые слова:} продовольственная безопасность,
агропродовольственная система, системные факторы, институциональная
среда, устойчивость, Республика Казахстан.

\begin{header}
АГРОӨНЕРКӘСІП ЖҮЙЕСІ ҚАУІПСІЗДІГІН ҚАЛЫПТАСТРУДАҒЫ ЖҮЙЕЛІК ФАКТОРЛАРЫ

\tsp{1}А.М.~Шілманова\envelope,
\tsp{1}Г.А.~Ибраева,
\tsp{1}А.С.~Шайнуров,
\tsp{2}Г.К.~Иляшова,
\tsp{3}С.А.~Юсупова,
\tsp{1}С.Ғ.~Тұңғышбай
\end{header}

\begin{affil}
\tsp{1}Қорқыт Ата атындағы Қызылорда университеті, Қызылорда, Казақстан,

\tsp{2}Әл-Фараби атындағы Қазақ ұлттық университеті, Алматы, Казақстан,

\tsp{3}Қазақ ұлттық хореография академиясы, Астана, Казақстан,

e-mail: altynajsilmanova1@gmail.com
\end{affil}

Жаһандық тұрақсыздықтың күшеюі және сыртқы әрі ішкі қатерлердің артуы
жағдайында азық-түлік қауіпсіздігі барған сайын жүйелік сипатқа ие
болып, оны талдаудың және қамтамасыз етудің дәстүрлі тәсілдерін қайта
қарауды талап етеді. Қазіргі агроазық-түлік жүйелері жоғары белгісіздік
жағдайында жұмыс істейді, бұл азық-түлік қауіпсіздігін тек өндіріс
көлемі мен азық-түліктің қолжетімділігі көрсеткіштері арқылы ғана емес,
сонымен қатар жүйенің орнықтылығы, келісімділігі және негізгі буындар
арасындағы өзара әрекеттесулердің қайта өндірілу қабілеті тұрғысынан
қарастырудың маңыздылығын арттырады.

Зерттеудің мақсаты агроазық-түлік жүйесінің азық-түлік қауіпсіздігін
қалыптастыратын жүйелік факторларды айқындау және олардың сыртқы және
ішкі күйзелістер жағдайында азық-түлікпен қамтамасыз етудің орнықтылығын
қамтамасыз етудегі рөлін талдау болып табылады. Зерттеудің әдіснамалық
негізін азық-түлік қауіпсіздігін агроазық-түлік жүйесінің тұтастай
қызмет етуінің нәтижесі ретінде қарастыруға мүмкіндік беретін жүйелік,
институционалдық және салыстырмалы талдау тәсілдері құрайды. Ақпараттық
база ретінде ФАО, ЭЫДҰ материалдары, ғылыми жарияланымдар және ресми
статистикалық деректер пайдаланылды.

Зерттеу нәтижесінде Қазақстан Республикасында азық-түлік қауіпсіздігін
қамтамасыз ету моделі негізінен жекелеген шаруашылық субъектілерін
қолдауға бағытталған қаржылық-инструмент\-алдық шараларға сүйенетіні
анықталды. Бұл ретте агроазық-түлік тізбектерінің қатысушыларын
үйлестіру, кооперациялау және орнықты интеграциялау тетіктері жеткілікті
деңгейде дамымаған және жүйе құраушы функция атқармайды. Ұзақ мерзімді
перспективада азық-түлікпен қамтамасыз етудің орнықтылығы бюджеттік
қолдау көлемінен гөрі институционалдық ортаның сапасына және
агроазық-түлік жүйесі буындары арасындағы өзара әрекеттесулердің қайта
өндірілуіне тәуелді екендігі жөнінде қорытынды жасалды.

{\bfseries Түйін сөздер:} азық-түлік қауіпсіздігі, агроазық-түлік жүйесі,
жүйелік факторлар, институционалдық орта, орнықтылық, Қазақстан
Республикасы.

\begin{header}
SYSTEMIC FACTORS IN THE FORMATION OF FOOD SECURITY WITHIN THE AGRI-FOOD SYSTEM

\tsp{1}A.M.~Shilmanova\envelope,
\tsp{1}G.А.~Ibraeva,
\tsp{1}A.S.~Shainurov,
\tsp{2}G.K.~Ilyashova,
\tsp{3}S.A.~Yusupova,
\tsp{1}S.G.~Tungyshbai
\end{header}

\begin{affil}
\tsp{1}Korkyt Ata Kyzylorda University, Kyzylorda, Kazakhstan,

\tsp{2}Al-Farabi Kazakh National University, Almaty, Kazakhstan,

\tsp{3}Kazakh National Academy of Choreography, Astana, Kazakhstan,

e-mail: altynajsilmanova1@gmail.com
\end{affil}

In the context of increasing global instability and the growing impact
of external and internal shocks, food security is increasingly acquiring
a systemic nature and requires a reconsideration of traditional
analytical approaches. Contemporary agri-food systems operate under
conditions of high uncertainty, which makes it insufficient to assess
food security solely through indicators of agricultural production and
food availability. Greater attention must be paid to the sustainability,
coordination, and reproducibility of interactions between the key
components of the agri-food system.

The purpose of this study is to identify and analyze the systemic
factors shaping food security in the agri-food system and to determine
their role in ensuring the sustainability of food provision under
conditions of external and internal shocks. The methodological framework
of the study is based on a combination of systemic, institutional, and
comparative approaches, which allow food security to be examined as an
outcome of the functioning of the agri-food system as a whole. The
information base includes materials from international organizations,
scientific publications, and official statistical data, including
sources from FAO, OECD, and national analytical reports.

The results of the study show that in the Republic of Kazakhstan the
food security model largely relies on financial and instrumental
measures of state support aimed at individual economic entities. At the
same time, mechanisms of coordination, cooperation, and sustainable
integration of participants in agri-food value chains remain
insufficiently developed and do not perform a system-forming function.
It is concluded that in the long term the sustainability of food
security depends less on the scale of budgetary support and more on the
quality of the institutional environment and the reproducibility of
interactions within the agri-food system.

{\bfseries Keywords:} food security, agri-food system, systemic factors,
institutional environment, sustainability, Republic of Kazakhstan.

\begin{multicols}{2}
{\bfseries Введение.} В условиях усиливающейся глобальной нестабильности
вопросы продовольственной безопасности приобретают выраженный системный
характер и выходят за рамки традиционного анализа, основанного
исключительно на объемах сельскохозяйственного производства и уровне
самообеспеченности. Нарушения логистических цепочек, рост цен на
продовольственные ресурсы, климатические и технологические риски, а
также высокая волатильность мировых рынков усиливают уязвимость
национальных агропродовольственных систем и актуализируют необходимость
пересмотра подходов к формированию продовольственной безопасности.

Современные международные исследования подчеркивают, что
продовольственная безопасность определяется не только физической
доступностью продовольствия, но и характером функционирования
агропродовольственной системы в целом, включая процессы переработки,
распределения, рыночного взаимодействия и координации участников. В этом
контексте она рассматривается как результат действия совокупности
взаимосвязанных факторов, формирующих устойчивость, адаптивность и
воспроизводимость агропродовольственных цепочек в условиях внешних и
внутренних воздействий.

Для стран с протяженной территорией, высокой региональной
дифференциацией и значительной долей малых и средних производителей
данная проблематика приобретает особую значимость. Агропродовольственная
система Республики Казахстан характеризуется сочетанием достаточной
сырьевой базы и одновременно повышенной чувствительности к
экономическим, логистическим и институциональным шокам, что во многом
обусловлено фрагментарностью связей между производственными,
перерабатывающими и распределительными звеньями.

Несмотря на наличие значительного массива научных исследований,
посвященных продовольственной безопасности, в литературе сохраняется
недостаточная разработанность системного подхода к анализу факторов ее
формирования. Большинство работ сосредоточено либо на макроэкономических
показателях и объемах производства, либо на отдельных сегментах
продовольственного рынка, тогда как комплексное рассмотрение
институциональных, организационных и координационных факторов остается
фрагментарным.

Целью настоящего исследования является выявление и анализ системных
факторов формирования продовольственной безопасности
агропродовольственной системы, а также определение их роли в обеспечении
устойчивости системы в условиях внешних и внутренних вызовов.

Гипотеза исследования заключается в том, что уровень продовольственной
безопасности определяется не только масштабами сельскохозяйственного
производства, но и характером действия системных факторов, связанных с
организацией, координацией и устойчивостью функционирования
агропродовольственной системы в целом.

Научная новизна исследования заключается в интерпретации
продовольственной безопасности не как результата действия отдельных
инструментов аграрной политики или объемов сельскохозяйственного
производства, а как следствия функционирования агропродовольственной
системы в целом, определяемого совокупностью системных факторов. В
работе обосновано, что ключевую роль в формировании устойчивости
продовольственного обеспечения играют институциональная архитектура,
степень согласованности и воспроизводимости взаимодействий между
производственными, перерабатывающими и сбытовыми звеньями, а также
наличие встроенных механизмов адаптации к внешним и внутренним шокам. В
отличие от преобладающих в научной литературе подходов, ориентированных
преимущественно на показатели доступности и самообеспеченности, в
исследовании предложен системный анализ продовольственной безопасности,
позволяющий выявить ограниченность финансово-инструментальной модели
поддержки и обосновать необходимость смещения акцентов государственной
аграрной политики в сторону развития устойчивых механизмов координации и
интеграции участников агропродовольственных цепочек.

Таким образом, в существующих исследованиях недостаточно раскрыта роль
системных факторов, обеспечивающих согласованность и воспроизводимость
агропродовольственной системы как целостного механизма формирования
продовольственной безопасности.

{\bfseries Литературный обзор.} Вопрос продовольственной безопасности в
последние десятилетия приобрел характер глобального приоритета и
рассматривается как ключевая задача социально-экономического развития
как развитых, так и развивающихся стран. При этом эмпирические
исследования показывают, что высокий уровень экономического развития и
наличие разветвленной системы аграрной поддержки не устраняют
структурные уязвимости агропродовольственных систем. Так, на примере
стран Европейского союза выявлено, что концентрация производства,
неравномерное распределение поддержки и ослабление позиций малых
производителей формируют долгосрочные риски для устойчивости
продовольственной безопасности, несмотря на развитую институциональную
среду {[}1{]}.

Дополнительные сравнительные оценки достижения Цели устойчивого развития
SDG-2 подтверждают, что даже среди экономически развитых стран
сохраняются существенные различия в уровне продовольственной
безопасности. Это указывает на ограниченность универсальных
стратегических решений и необходимость учета совокупности структурных и
институциональных факторов, определяющих характер функционирования
агропродовольственных систем. Данный вывод находит подтверждение в
исследованиях, подчеркивающих ключевую роль координации и
согласованности действий участников системы в преодолении структурных
ограничений продовольственной безопасности {[}2{]}.

На национальном уровне большинство стран разрабатывают
специализированные стратегии, программы и дорожные карты в сфере
продовольственной безопасности, которые, как правило, выходят за рамки
узкого аграрного регулирования и охватывают вопросы институционального
развития, координации политики, поддержки фермерских хозяйств, развития
кооперации и интеграции участников агропродовольственных цепочек. Вместе
с тем в научных исследованиях подчеркивается, что наличие формальных
стратегий и программ само по себе не гарантирует достижения устойчивых
результатов, если механизмы их реализации характеризуются
фрагментарностью и недостаточной согласованностью.

Ряд исследований показывает, что отсутствие согласованности между
уровнями управления существенно снижает результативность политики в
сфере продовольственной безопасности. Анализ многоуровневого управления
аграрными системами свидетельствует о том, что универсальные
регуляторные инструменты, реализуемые на национальном уровне, нередко
вступают в противоречие с локальными институциональными практиками. Это
ограничивает адаптационный потенциал производителей и препятствует
устойчивому функционированию агропродовольственных систем {[}3{]}.

В этой связи в современной литературе усиливается интерес к системному
анализу факторов формирования продовольственной безопасности, включая
характер взаимодействия между участниками агропродовольственной системы,
формы координации и согласования интересов, а также качество
институциональной среды. Такой подход позволяет рассматривать
продовольственную безопасность не только как результат производства и
распределения ресурсов, но и как следствие функционирования
агропродовольственной системы в целом.

Современные исследования все чаще трактуют продовольственную
безопасность как результат функционирования агропродовольственной
системы, в которой ключевую роль играют не только объемы
сельскохозяйственного производства, но и системные факторы, определяющие
согласованность, устойчивость и воспроизводимость агропродовольственных
процессов. Данный подход отражает переход от узкого понимания
продовольственной безопасности как физической доступности продовольствия
к системному анализу устойчивости и координации на различных уровнях
управления.

Ряд работ подчеркивает, что влияние институциональной среды на
показатели продовольственной безопасности носит опосредованный характер
и во многом определяется согласованностью действий и механизмов
управления в условиях кризисов. Анализ продовольственной безопасности в
странах Западных Балкан показывает, что фрагментарность политики и
недостаточная координация между уровнями управления усиливают уязвимость
агропродовольственных систем и ограничивают результативность мер
государственной поддержки даже при наличии формальных стратегий и
программ {[}4{]}.

Существенное внимание в литературе уделяется анализу форм взаимодействия
и координации экономических агентов в агропродовольственных цепочках.
Исследование аграрных социальных сервисов в Китае демонстрирует, что
специализированные сервисные организации и кооперативные структуры
выступают важным системным фактором повышения устойчивости
продовольственного производства, способствуя интеграции мелких
производителей в современные производственные и рыночные процессы и
повышая эффективность использования ресурсов {[}5{]}.

Отдельное направление исследований связано с анализом распределения
выгод и рисков между участниками агропродовольственной системы. В этом
контексте категория справедливости (fairness) рассматривается как один
из системных факторов устойчивости продовольственных систем.
Подчеркивается, что асимметрия рыночной власти, информационные
дисбалансы и исключение уязвимых групп производителей из цепочек
создания стоимости подрывают потенциал продовольственной безопасности
даже при наличии формальных регуляторных механизмов {[}6{]}.

В казахстанской научной литературе продовольственная безопасность
рассматривается преимущественно в контексте структурных и
институциональных ограничений развития агропродовольственной системы.
Авторы отмечают, что при наличии значительного сельскохозяйственного
потенциала сохраняется высокая зависимость внутреннего рынка от импорта
по ряду ключевых товарных позиций, а продовольственная безопасность
формируется в условиях выраженной региональной неоднородности,
обусловленной различиями в производственном потенциале, логистической
обеспеченности и уровне социально-экономического развития территорий
{[}7--9{]}. Эти факторы усиливают уязвимость агропродовольственной
системы к внешним шокам и ценовым колебаниям.

Значительное место в отечественных исследованиях занимает анализ
экономической доступности продовольствия. Высокая доля расходов
домохозяйств на продукты питания в ряде регионов свидетельствует о
сохранении рисков продовольственной уязвимости населения. Подтверждается
действие закона Энгеля, в соответствии с которым рост доходов
сопровождается снижением доли расходов на продовольствие, что позволяет
рассматривать структуру потребления как важный индикатор
продовольственной безопасности и социально-экономического развития
{[}9{]}. Одновременно выявляется региональное неравенство, указывающее
на недостаточную согласованность аграрной, социальной и региональной
политики {[}10--11{]}.

Методологические аспекты оценки продовольственной безопасности также
остаются предметом активной научной дискуссии. Анализ используемых
индикаторов показывает, что в большинстве исследований доминируют
показатели доступности и стоимости продовольствия, тогда как
институциональные, организационные и экологические параметры часто
остаются вне поля системного анализа. Это ограничивает возможности
комплексной оценки системных факторов формирования продовольственной
безопасности {[}12{]}.

Анализ научных источников показывает, что продовольственная безопасность
формируется под воздействием совокупности системных факторов,
определяющих характер функционирования агропродовольственной системы.
Несмотря на значительный массив исследований, в литературе сохраняется
фрагментарность в анализе взаимосвязи между институциональной средой,
формами взаимодействия участников и результатами продовольственной
безопасности, что обуславливает необходимость дальнейшего системного
анализа применительно к национальным условиям.

{\bfseries Материалы и методы.} В рамках настоящего исследования
использован комплексный методологический подход, основанный на сочетании
системного анализа, институционального анализа и сравнительного анализа.
Такой выбор методов обусловлен пониманием продовольственной безопасности
как результата действия совокупности системных факторов, формируемых в
процессе функционирования агропродовольственной системы под воздействием
экономических, институциональных и организационно-координационных
условий.

Информационную базу исследования составили материалы международных и
национальных организаций, научные публикации, статистические данные и
официальные документы, отражающие различные аспекты формирования
продовольственной безопасности. В частности, использованы аналитические
и статистические материалы ФАО, Всемирного банка и ОЭСР, а также
результаты эмпирических и обзорных исследований, опубликованных в
рецензируемых научных журналах. Для анализа национальной специфики
продовольственной безопасности Республики Казахстан привлечены данные
официальной статистики, отраслевые обзоры, аналитические отчеты и
публикации отечественных исследователей.

Методы исследования. Ключевым методологическим инструментом исследования
выступает системный анализ, позволяющий рассматривать продовольственную
безопасность как результат взаимодействия взаимосвязанных элементов
агропродовольственной системы, включая производство, переработку,
распределение, институциональную среду и механизмы координации
участников. Использование системного подхода обеспечивает возможность
выявления не отдельных факторов, а их совокупного влияния на
устойчивость и воспроизводимость продовольственного обеспечения.

Институциональный анализ применен для оценки роли формальных и
неформальных институтов, механизмов координации и организационных
условий, которые выступают системными факторами формирования
продовольственной безопасности. В рамках данного метода внимание уделено
не отдельным инструментам государственной поддержки, а характеру
институциональной архитектуры и степени согласованности действий
участников агропродовольственной системы.

Сравнительный анализ использован для сопоставления международных и
национальных подходов к обеспечению продовольственной безопасности, а
также для выявления общих и специфических закономерностей
функционирования агропродовольственных систем в различных
институциональных и экономических условиях. Это позволило рассмотреть
системные факторы формирования продовольственной безопасности в более
широком контексте и избежать узконациональной интерпретации результатов.

В целях систематизации научных подходов и выявления ключевых направлений
исследовательской дискуссии применен контент-анализ научных публикаций,
посвященных проблемам продовольственной безопасности, продовольственного
суверенитета и устойчивости агропродовольственных систем. Особое
внимание уделялось работам, рассматривающим вопросы координации,
распределения выгод и рисков, а также организационных и
институциональных условий функционирования агропродовольственных
цепочек.

Анализ продовольственной безопасности в рамках исследования базируется
на рассмотрении совокупности показателей физической и экономической
доступности продовольствия, а также характеристик устойчивости
агропродовольственных систем. При этом акцент сделан на интерпретации
данных показателей через призму системных факторов, включая
согласованность политики, уровень координации между участниками
агропродовольственных цепочек и организационные условия воспроизводства
продовольственного предложения. Такой подход позволяет выйти за рамки
узкого анализа производственных показателей и рассматривать
продовольственную безопасность как результат функционирования
агропродовольственной системы в целом.

Следует отметить, что результаты исследования основаны преимущественно
на вторичных данных и опубликованных научных источниках, что
ограничивает возможности анализа неформальных институциональных практик
и локальных механизмов взаимодействия. Вместе с тем использование
междисциплинарного, системного и сравнительного подходов обеспечивает
достаточную обобщаемость выводов и формирует методологическую основу для
последующих эмпирических исследований, ориентированных на углубленный
анализ системных факторов формирования продовольственной безопасности.

{\bfseries Обсуждение и результаты.} В Республике Казахстан государственная
поддержка агропромышленного комплекса формируется в рамках
многоуровневой институциональной архитектуры, ориентированной на
финансово-экономическое стимулирование предпринимательской деятельности
в сельском хозяйстве. Базовыми инструментами такой поддержки выступают
льготное налогообложение, субсидирование инвестиционных и текущих
затрат, а также льготное кредитование сезонных сельскохозяйственных
работ, что позволяет снизить финансовую нагрузку на субъекты АПК и
обеспечить относительную стабильность производственного процесса.

Координация и реализация ключевых мер государственной поддержки
осуществляется через систему квазигосударственных институтов развития,
центральное место среди которых занимает Национальный управляющий
холдинг Байтерек. В структуре холдинга сосредоточены специализированные
организации, обеспечивающие доступ субъектов АПК к финансовым ресурсам,
инвестициям и инструментам обновления материально-технической базы. К
числу таких институтов относятся Аграрная кредитная корпорация,
осуществляющая льготное кредитование и финансирование агробизнеса через
различные финансовые каналы, а также КазАгроФинанс, специализирующаяся
на лизинге сельскохозяйственной техники и оборудования.

Функциональная направленность указанных институтов свидетельствует о
преобладании финансово-инструментального подхода к развитию
агропромышленного комплекса, при котором основной акцент делается на
расширение доступа субъектов агробизнеса к финансовым ресурсам,
технологическую модернизацию и повышение производительности труда. При
этом поддержка преимущественно ориентирована на отдельные хозяйствующие
субъекты, тогда как формирование устойчивой связанности между
производственными, перерабатывающими и сбытовыми звеньями
агропродовольственной системы рассматривается как вторичный результат, а
не как самостоятельное направление государственной политики.

В результате сложившаяся модель государственной поддержки АПК создает
базовые условия для сохранения производственного потенциала отрасли,
однако в ограниченной степени способствует развитию системных факторов
устойчивости агропродовольственной системы. Это формирует зависимость
продовольственной безопасности от масштабов бюджетной поддержки и
снижает роль встроенных механизмов взаимодействия и воспроизводства
связей между звеньями агропродовольственных цепочек в долгосрочной
перспективе.
\end{multicols}

\tcap{Таблица 1 - Институциональная архитектура государственной
поддержки агропромышленного комплекса Республики Казахстан}
\begin{longtblr}[
  label = none,
  entry = none,
]{
  width = \linewidth,
  colspec = {Q[10]Q[150]Q[249]Q[138]Q[134]Q[227]},
  cells = {font = \footnotesize},
  row{1} = {c},
  column{1} = {c},
  row{7} = {font=\itshape\footnotesize},
  cell{1}{2} = {font=\bfseries\footnotesize},
  cell{1}{3} = {font=\bfseries\footnotesize},
  cell{1}{4} = {font=\bfseries\footnotesize},
  cell{1}{5} = {font=\bfseries\footnotesize},
  cell{7}{1} = {c=6}{0.936\linewidth},
  hlines,
  vlines,
}
№ & Институт/ организация & Основные инструменты поддержки & Целевая группа получателей & Направлен\-ность воздействия & {\textbf{Потенциальный вклад в обеспечение}\\\textbf{продовольственной безопасности}} \\
1 & Министерство сельского хозяйства Республики Казахстан & Субсидирование инвестиционных затрат; субсидирование процентных ставок по кредитам и лизингу; льготное кредитование весенне-полевых и уборочных работ & Субъекты АПК, фермерские и крестьянские хозяйства, агропредприятия & Производство, финансирование & Стабилизация сельскохозяйственного производства и снижение финансовых рисков аграрного сектора \\
2 & Национальный управляющий холдинг Байтерек & Координация и реализация финансовых инструментов поддержки через дочерние организации & Субъекты АПК, малый и средний агробизнес & Финансирова\-ние, институциональная координация & Формирование централизованной системы доступа к финансовым ресурсам и повышение устойчивости агропродовольственной системы \\
3 & Аграрная кредитная корпорация & Льготное кредитование; субсидирование ставок вознаграждения; финансирование через кредитные товарищества, МФО и банки второго уровня & Субъекты АПК, МСП, фермерские хозяйства & Финансирова\-ние, производство & Расширение охвата агробизнеса финансовыми услугами и поддержание непрерывности производственного цикла \\
4 & КазАгроФинанс & Лизинг сельскохозяйственной техники и оборудования; кредитование & Субъекты АПК, сельхозпроизводители & Техническое оснащение, производство & Обновление материально-технической базы и повышение производительности сельского хозяйства \\
5 & Финансовые институты-партнеры (банки второго уровня, МФО, лизинговые компании) & Со-финансирование, кредитование и реализация льготных программ через институциональные каналы & Малый и средний агробизнес, фермерские хозяйства & Финансирова\-ние & Повышение доступности финансовых ресурсов и снижение институциональных барьеров входа в отрасль \\
\textit{Примечание - составлено авторами на основе анализа нормативно-правовых актов, официальных материалов Министерства сельского хозяйства Республики Казахстан и данных квазигосударственных институтов развития.} & & & & & 
\end{longtblr}

\begin{multicols}{2}
В условиях преобладания финансово-инструментального подхода особое
значение приобретает анализ механизмов взаимодействия и интеграции как
одного из системных факторов, дополняющих меры финансовой поддержки и
влияющих на устойчивость агропродовольственной системы. В отличие от
поддержки, ориентированной преимущественно на отдельные хозяйствующие
субъекты, такие механизмы связаны с формированием устойчивой связанности
между производственными, перерабатывающими, логистическими и сбытовыми
звеньями, а также с вовлечением финансовых и научных организаций в
единое пространство взаимодействия.

Механизмы взаимодействия в агропродовольственной системе проявляются в
устойчивых формах согласования интересов экономических агентов на
различных стадиях создания продовольственной стоимости. В отличие от
разрозненных рыночных трансакций, они опираются на повторяемость
взаимодействий, распределение рисков и выгод и наличие формальных либо
неформальных правил, обеспечивающих согласованность действий участников.
В условиях высокой неопределенности и уязвимости агропродовольственных
цепочек такие формы взаимодействия выступают не только источником
повышения эффективности, но и важным системным фактором устойчивости
продовольственной безопасности.

Исследования, посвященные анализу взаимодействия участников
агропродовольственной системы Республики Казахстан, показывают, что на
практике такие формы взаимодействия характеризуются фрагментарностью и
низкой устойчивостью {[}13{]}. Несмотря на наличие значительного числа
сельскохозяйственных производителей, прежде всего малых и средних
хозяйств, их включенность в процессы переработки и сбыта продукции
остается ограниченной. Существующие формы взаимодействия, как правило,
носят ситуативный характер и зависят от доступа к государственным мерам
поддержки либо отдельных рыночных возможностей, не формируя устойчивой
связанности между звеньями агропродовольственных цепочек.

С позиций обеспечения продовольственной безопасности данная особенность
имеет принципиальное значение. Недостаточная связанность и слабая
воспроизводимость взаимодействий между звеньями агропродовольственной
системы приводят к разрывам в цепочках создания продовольственной
стоимости, росту трансакционных издержек и снижению адаптационного
потенциала системы. Малые производители оказываются ограниченно
вовлеченными в переработку и распределение продукции, что сдерживает
устойчивость продовольственного предложения и усиливает зависимость
внутреннего рынка от внешних поставок по ряду товарных позиций.

Кооперативные и сетевые формы взаимодействия, рассматриваемые в научной
литературе как важный элемент устойчивых агропродовольственных систем, в
казахстанских условиях преимущественно выполняют вспомогательную роль и
не встроены в систему обеспечения продовольственной безопасности как
самостоятельный системный компонент. В отличие от моделей, в которых
такие формы взаимодействия используются для стабилизации доходов
производителей, сглаживания региональных диспропорций и обеспечения
непрерывности поставок, в Республике Казахстан они чаще проявляются как
производный элемент финансовой поддержки, а не как инструмент
долгосрочной организации и согласования агропродовольственных процессов.

В результате элементы согласования и интеграции в агропродовольственной
системе Республики Казахстан в настоящее время не выполняют
системообразующую функцию в формировании продовольственной безопасности.
Их ограниченная встроенность в общую логику функционирования системы
снижает устойчивость агропродовольственных цепочек и усиливает
уязвимость продовольственного обеспечения к внешним и внутренним шокам.
В этой связи формирование устойчивой связанности и воспроизводимости
взаимодействий между ключевыми звеньями агропродовольственной системы
следует рассматривать как одно из приоритетных направлений повышения
устойчивости продовольственной безопасности в долгосрочной перспективе.

Выявленные ограничения актуализируют необходимость обращения к
зарубежным концептуальным моделям, в которых устойчивость
продовольственного обеспечения трактуется как результат согласованного
функционирования агропродовольственной системы в целом. В международных
исследованиях все более широкое распространение получает системный
подход, рассматривающий продовольственную безопасность не как функцию
отдельных отраслей или инструментов поддержки, а как следствие качества
организации, связанности и согласованности процессов производства,
переработки, распределения и потребления продовольствия. В этом
контексте особый интерес представляет концептуальная рамка устойчивости
агропродовольственных систем, предложенная ФАО, позволяющая
структурировать ключевые системные характеристики устойчивости в
условиях внешних и внутренних шоков (рисунок 1).
\end{multicols}

\fig{e/image5}[Рис.1 - Концептуальная рамка анализа устойчивости
агропродовольственных систем\\\normalfont{\emph{Примечание -- источник {[}14{]}}}]

\begin{multicols}{2}
В соответствии с данной концептуальной рамкой устойчивость
агропродовольственной системы интерпретируется как ее способность
сохранять функциональность и восстанавливаться после воздействия шоков
за счет сочетания устойчивости к возмущениям, способности к поглощению
негативных эффектов, скорости восстановления и адаптивности. Как
показано на рисунке 2, уровень продовольственной безопасности во времени
определяется не столько масштабами самого шока, сколько состоянием
системных факторов, обеспечивающих избыточность связей, гибкость цепочек
поставок и согласованность функционирования агропродовольственной
системы.

В рамках данной модели именно совокупность взаимосвязанных элементов
системы определяет глубину снижения продовольственной безопасности в
период кризисных воздействий и скорость последующего восстановления.
Недостаточная развитость таких системных характеристик, напротив,
приводит к более выраженному снижению устойчивости и усиливает
зависимость продовольственного обеспечения от внешних факторов и
экстренных мер государственной поддержки.

Сопоставление представленных концептуальных моделей устойчивости
агропродовольственной системы с практикой Республики Казахстан
показывает, что ключевые системные характеристики, определяющие
устойчивость системы - избыточность связей, гибкость цепочек поставок и
адаптивность взаимодействия между звеньями, в национальной модели
выражены в ограниченной степени. При доминировании
финансово-инструментальных мер государственной поддержки устойчивость
продовольственного обеспечения в большей мере зависит от масштабов
бюджетного финансирования и административных решений, чем от встроенных
механизмов воспроизводства связей внутри агропродовольственной системы,
что соответствует траектории менее устойчивой системы, представленной на
концептуальной кривой.

В целом результаты исследования позволяют рассматривать характер
организации агропродовольственной системы, модель государственной
поддержки и степень связанности ее ключевых звеньев как системные
факторы формирования продовольственной безопасности, определяющие
устойчивость и адаптационный потенциал системы в условиях внешних и
внутренних шоков.
\end{multicols}

\fig[0.6\textwidth]{e/image6}[Рис.2 - Устойчивость продовольственной системы как способность
обеспечивать продовольственную безопасность во времени в условиях
внешних возмущений\\\normalfont{\emph{Примечание -- источник {[}15{]}}}]

\begin{multicols}{2}
{\bfseries Выводы.} Проведенное исследование показало, что
продовольственная безопасность в современных условиях формируется не как
прямое следствие объемов сельскохозяйственного производства или
масштабов государственной поддержки, а как результат функционирования
агропродовольственной системы в целом, определяемого совокупностью
системных факторов. К числу таких факторов относятся характер
институциональной архитектуры, степень согласованности и
воспроизводимости взаимодействий между ключевыми звеньями
агропродовольственных цепочек, а также наличие встроенных механизмов
адаптации к внешним и внутренним шокам.

Установлено, что в Республике Казахстан модель обеспечения
продовольственной безопасности в значительной степени опирается на
финансово-инструментальный подход, ориентированный на поддержку
отдельных хозяйствующих субъектов. Данная модель позволяет поддерживать
производственный потенциал агропромышленного комплекса и снижать
краткосрочные финансовые риски, однако в ограниченной степени
способствует формированию устойчивой связанности между
производственными, перерабатывающими и сбытовыми звеньями
агропродовольственной системы. В результате устойчивость
продовольственного обеспечения во многом зависит от масштабов бюджетной
поддержки и административных решений.

Показано, что механизмы кооперации, сетевого взаимодействия и
партнерских связей в агропродовольственной системе Республики Казахстан
носят фрагментарный характер и не выполняют системообразующей функции.
Ограниченная интеграция малых и средних производителей в процессы
переработки и распределения продукции снижает адаптационный потенциал
агропродовольственной системы и усиливает ее уязвимость к ценовым,
логистическим и институциональным шокам.

На основе полученных результатов авторами предложена обобщающая схема
формирования продовольственной безопасности агропродовольственной
системы, отражающая взаимосвязь институциональной архитектуры, форм
взаимодействия участников и характеристик устойчивости системы (рисунок
3).
\end{multicols}

{\bfseries Рис.3 - Схема формирования устойчивости продовольственной
безопасности агропродовольственной системы}

\emph{Примечание - составлено авторами}

\begin{multicols}{2}
Сопоставление национальной практики с концептуальной рамкой устойчивости
агропродовольственных систем ФАО позволило установить, что ключевые
характеристики устойчивости - избыточность связей, гибкость цепочек
поставок и адаптивность взаимодействий - выражены в национальной модели
в ограниченной степени, что обуславливает зависимость восстановления
продовольственной безопасности в кризисные периоды от экстренных мер
государственной поддержки.

С практической точки зрения полученные результаты указывают на
необходимость смещения акцентов государственной аграрной политики от
преимущественно финансово-инструментальной поддержки отдельных
хозяйствующих субъектов к формированию устойчивых механизмов координации
и интеграции участников агропродовольственных цепочек. Речь идет не о
замещении действующих мер поддержки, а о их институциональном дополнении
за счет развития кооперативных, сетевых и партнерских форм
взаимодействия между производственными, перерабатывающими,
логистическими и сбытовыми звеньями. Усиление таких механизмов позволяет
повысить воспроизводимость связей внутри агропродовольственной системы,
снизить трансакционные издержки и повысить адаптационный потенциал
продовольственного обеспечения в условиях внешних и внутренних шоков,
что имеет принципиальное значение для долгосрочной устойчивости
продовольственной безопасности Республики Казахстан.
\end{multicols}

\begin{center}
{\bfseries Литература}
\end{center}

\begin{refs}
1. Moreno-Pérez, O. M., Arnalte-Mur, L., Cerrada-Serra, P.,
Sumpsi-Viñas, J., \& Ortiz-Miranda, D. Actions to strengthen the
contribution of small farms and small food businesses to food security
in Europe// Food Security. - 2024.- Vol.16(1). - P.243--259.
\href{https://doi.org/10.1007/s12571-023-01421-0}{DOI
10.1007/s12571-023-01421-0}.

2. Smyth S.J., Webb S.R., Phillips P.W.B. The role of public-private
partnerships in improving global food security//Global Food Security. -
2021. -Vol 31:100588. DOI 10.1016/j.gfs.2021.100588.

3. Sidibé A., Totin E., Thompson-Hall M., Obisis L., Traoré P.C.
Multi-scale governance in agricultural systems: Interplay between
national and local institutions around the production dimension of food
security in Mali // NJAS - Wageningen Journal of Life Sciences.
-2018.-Vol.84.- P.94-102. DOI 10.1016/j.njas.2017.09.001.

4. Jambor A., Varga Á. Food security and crises: evidence from the
Western Balkans //
\href{https://www.researchgate.net/journal/Agriculture-Food-Security-2048-7010?_tp=eyJjb250ZXh0Ijp7ImZpcnN0UGFnZSI6InB1YmxpY2F0aW9uIiwicGFnZSI6InB1YmxpY2F0aW9uIn19}{Agriculture
\& Food Security}.- 2025.-Vol.13(1). DOI 10.1007/s40066-024-00514-z.

5. Cai B., Wang L., Shi F., Abate M. C., Geremew B., \& Addisu, A. K.
(2024). Agricultural socialized services and Chinese food security:
Examining the threshold effect of land tenure change//Frontiers in
Sustainable Food Systems. - 2024.-Vol.8:1371520.
\href{https://doi.org/10.3389/fsufs.2024.1371520}{DOI
10.3389/fsufs.2024.1371520}.

6. Onyeaka H., et al. Achieving fairness in the food system// Food and
Energy Security. - 2024.- Vol.13(4): e572. DOI 10.1002/fes3.572.

7. Дуламбаева Р., Жунусова А. Оценка продовольственной безопасности
Казахстана по компонентам доступности // Мемлекеттік басқару және
мемлекеттік қызмет. - 2024. - № 1(88). - С.175-188. DOI
10.52123/1994-2370-2024-1176.

8. Дедкова Е.Г., Коростелкина И.А., Жидова О.П. Актуальные аспекты
обеспечения продовольственной безопасности // Journal of Economic
Research \& Business Administration. - 2019. -№ 3 (129). - С.172-187.
DOI 10.26577/be-2019-3-e15.

9. Беделбаева А.Е., Шарипов А.К., Досумова Ж.С. ЕАЭО жағдайында
Қазақстанның азық-түлік қауіпсіздігін қамтамасыз ету // Аграрлық нарық
проблемалары. -2021. - № 3. - С.24-30 DOI 10.46666/2021-3.2708-9991.02.

10. Jumabayev S., Dulambayeva R., Kussainova L., Yesmagambetov D.
Approaches to assessing the food security of the regions of Kazakhstan
in modern conditions // Public Policy and Administration. - 2023.-Vol.
22(4) - P.503-518.

11. Нурланова Н.К., Альжанова Ф.Г., Днишев Ф.М. (2023). Методы и
практика оценки уровня инклюзивного регионального развития Казахстана.
Экономика: стратегия и практика. - № 18(4). - С.109-126, DOI
10.51176/1997-9967-2023-4-109-126

12. Cadavid, L., Arulnathan, V., Pelletier, N. (2024). Food security and
food sovereignty: A review of commonly used indicators and consideration
of environmental sustainability aspects//Sustainability. -2024.-Vol.
16(24): 11034. \href{https://doi.org/10.3390/su162411034}{DOI
10.3390/su162411034}

13. Юсупова С., и др. «Формирование предпринимательских сетей в
агропромышленном комплексе Казахстана: институциональные и
организационные предпосылки»//Вестник КазУТБ. -2025. -№ 2(27). - С.
571-584. DOI 10.58805/kazutb.v.2.27-1011.

14. FAO.2021. In Brief the State of Food and Agriculture 2021. Making
agrifood systems more resilient to shocks and stresses// Rome, Italy. -
2021. - 28p. \href{https://doi.org/10.4060/cb7351en}{DOI
10.4060/cb7351en}.

15. Tendall, D. M., Joerin, J., Kopainsky, B., Edwards, P., Shreck, A.,
Le, Q. B., Kruetli, P., \& Grant, M. (2015). Food system resilience:
Defining the concept//Global Food Security. -2015.-Vol.6.- P.17-23. DOI
10.1016/j.gfs.2015.08.001.
\end{refs}

\begin{center}
{\bfseries References}
\end{center}

\begin{refs}
1. Moreno-Pérez, O. M., Arnalte-Mur, L., Cerrada-Serra, P.,
Sumpsi-Viñas, J., \& Ortiz-Miranda, D. Actions to strengthen the
contribution of small farms and small food businesses to food security
in Europe// Food Security. - 2024.- Vol.16(1). - P.243--259.
\href{https://doi.org/10.1007/s12571-023-01421-0}{DOI
10.1007/s12571-023-01421-0}.

2. Smyth S.J., Webb S.R., Phillips P.W.B. The role of public-private
partnerships in improving global food security//Global Food Security. -
2021. -Vol 31:100588. DOI 10.1016/j.gfs.2021.100588.

3. Sidibé A., Totin E., Thompson-Hall M., Obisis L., Traoré P.C.
Multi-scale governance in agricultural systems: Interplay between
national and local institutions around the production dimension of food
security in Mali // NJAS - Wageningen Journal of Life Sciences.
-2018.-Vol.84.- P.94-102. DOI 10.1016/j.njas.2017.09.001.

4. Jambor A., Varga Á. Food security and crises: evidence from the
Western Balkans //
\href{https://www.researchgate.net/journal/Agriculture-Food-Security-2048-7010?_tp=eyJjb250ZXh0Ijp7ImZpcnN0UGFnZSI6InB1YmxpY2F0aW9uIiwicGFnZSI6InB1YmxpY2F0aW9uIn19}{Agriculture
\& Food Security}.- 2025.-Vol.13(1). DOI 10.1007/s40066-024-00514-z.

5. Cai B., Wang L., Shi F., Abate M. C., Geremew B., \& Addisu, A. K.
(2024). Agricultural socialized services and Chinese food security:
Examining the threshold effect of land tenure change//Frontiers in
Sustainable Food Systems. - 2024.-Vol.8:1371520.
\href{https://doi.org/10.3389/fsufs.2024.1371520}{DOI
10.3389/fsufs.2024.1371520}.

6. Onyeaka H., et al. Achieving fairness in the food system// Food and
Energy Security. - 2024.- Vol.13(4): e572. DOI 10.1002/fes3.57.

7. Dulambaeva R., Zhunusova A. Ocenka prodovol' stvennoj
bezopasnosti Kazahstana po komponentam dostupnosti // Memlekettіk
basқaru zhәne memlekettіk қyzmet. - 2024. - № 1(88). - S.175-188. DOI
10.52123/1994-2370-2024-1176. {[}in Russian{]}

8. Dedkova E.G., Korostelkina I.A., Zhidova O.P.
Aktual' nye aspekty obespechenija
prodovol' stvennoj bezopasnosti // Journal of Economic
Research \& Business Administration. - 2019. -№ 3 (129). - S.172-187.
DOI 10.26577/be-2019-3-e15. {[}in Russian{]}

9. Bedelbaeva A.E., Sharipov A.K., Dosumova Zh.S. EAJeO zhaғdajynda
Қazaқstannyң azyқ-tүlіk қauіpsіzdіgіn қamtamasyz etu // Agrarlyқ naryқ
problemalary. -2021. - № 3. - S.24-30
DOI 10.46666/2021-3.2708-9991.02. {[}in Kazakh{]}
\end{refs}

\begin{info}
\hspace{1em}\emph{{\bfseries Сведения об авторах}}

Шильманова А.М.- кандидат экономических наук, Кызылординский
университет им. Коркыт Ата, Кызыорда, Казахстан, е-mail:
altynajsilmanova1@gmail.com;

Ибраева Г.А.- магистр экономических наук, Кызылординский университет
им.  Коркыт Ата, Кызыорда, Казахстан, е-mail:
ibraevagulajym@gmail.com;

Шайнуров А.С.- кандидат экономических наук, Кызылординский университет
им. Коркыт Ата, Кызыорда, Казахстан, е-mail: kaup@mail.ru;

Иляшова Г.Ж.- PhD, Казахский национальный университет им. аль-Фараби
г.  Алматы, Казахстан, е-mail: ilyashova.g@kaznu.kz;

Юсупова С.А.- кандидат экономических наук, Казахской национальной
академии хореографии, е-mail: Abay.saltanat@bk.ru;

Тұңғышбай С.Ғ.- магистр технологических наук, докторант,
Кызылординский университет им.- Коркыт Ата, Кызылорда, е-mail:
macqara@gmail.com.

\hspace{1em}\emph{{\bfseries Information~about~the~authors}}

Shilmanova A.M.- candidate of economic sciences, Korkyt Ata Kyzylorda
University, Kyzylorda, Kazakhstan, е-mail: a.shilmanova@bk.ru;

Ibraeva G.- master of Economic Sciences, Korkyt Ata Kyzylorda
University, Kyzylorda, Kazakhstan, e-mail: ibraevagulajym@gmail.com;

Shainurov A.S.- candidate of economic sciences, Korkyt Ata Kyzylorda
University, Kyzylorda, Kazakhstan, e-mail: kaup@mail.ru;

Ilyashova G. K.- PhD, Al-Farabi Kazakh National University, Almaty,
Kazakhstan, е-mail: ilyashova.g@kaznu.kz;

Yusupova S. A.- candidate of economic sciences, Kazakh National
Academy of Choreography, е-mail: Abay.saltanat@bk.ru;

Tungyshbai S. G.- master of Technological Sciences, doctoral student,
Korkyt Ata Kyzylorda University, Kyzylorda, Kazakhstan, e-mail:
macqara@gmail.com.
\end{info}
